\documentclass[../../main.tex]{subfiles}
\graphicspath{{../../image/}{../../figure/}}
\begin{document}

\section{标量守恒律的一个松弛格式}\label{section:ch17-标量松弛格式}

这里的目标是借助补偿紧致性理论，在方程组\eqref{align:ch05-relaxed-system}中令 $\mu\downarrow0$ 取极限。
这样一个练习可服务于双重目的：当人们对(5.2.18)本身作为某种物理过程的模型感兴趣时，它将展示松弛到由标量守恒律\eqref{align:ch17-1741}所支配的局部平衡；作为副产品，它将确立\eqref{align:ch17-1741}的 Cauchy 问题解的存在性，并提示一种计算它们的方法。为后一目的，宜使未松弛方程组(5.2.18)尽可能简单，即半线性，
\begin{align}
\begin{cases}
\partial_t u(x,t)+\partial_x v(x,t)=0\\
\partial_t v(x,t)+a^2\partial_x u(x,t)+\dfrac{1}{\mu}[v(x,t)-f(u(x,t))]=0,
\end{cases}
\label{align:ch17-1751}
\end{align}
其中 $a$ 为某正常数。为简化分析，我们在此只处理这个半线性方程组。此时读者最好回顾\cref{section:ch05-相对熵与经典解的稳定性}中给出的松弛理论导引。

设 $f'$ 有界，并选取 $a$ 足够大，使严格次特征条件（见\eqref{align:ch05-subcharacteristic}）
\begin{align}
-a+\delta<f'(u)<a-\delta,\qquad u\in(-\infty,\infty),
\label{align:ch17-1752}
\end{align}
对某 $\delta>0$ 成立。我们通过假定 $f(0)=0$ 来规范化 $v$。

\eqref{align:ch17-1751}的熵--熵流对 $(\eta(u,v),q(u,v))$ 满足线性双曲组
\begin{align}
\begin{cases}
q_u(u,v)-a^2\eta_v(u,v)=0\\
q_v(u,v)-\eta_u(u,v)=0,
\end{cases}
\label{align:ch17-1753}
\end{align}
其通解为
\begin{align}
\begin{cases}
\eta(u,v)=r(au+v)+s(au-v)\\
q(u,v)=ar(au+v)-as(au-v).
\end{cases}
\label{align:ch17-1754}
\end{align}

次特征条件\eqref{align:ch17-1752}蕴含曲线 $v=f(u)$ 对\eqref{align:ch17-1753}而言处处非特征；因此，对标量守恒律\eqref{align:ch17-1741}给定任意熵--熵流对 $(\hat\eta(u),\hat q(u))$，就可构造\eqref{align:ch17-1751}的熵--熵流对 $(\eta(u,v),q(u,v))$，使其满足 Cauchy 数据
\begin{align}
\eta(u,f(u))=\hat\eta(u),\qquad q(u,f(u))=\hat q(u),\qquad u\in(-\infty,\infty).
\label{align:ch17-1755}
\end{align}
对 $u$ 微分\eqref{align:ch17-1755}，并利用 $\hat q'(u)=\hat\eta'(u)f'(u)$，连同\eqref{align:ch17-1753}与\eqref{align:ch17-1752}，推出 $\eta_v(u,f(u))=0$。将此与\eqref{align:ch17-1755}、\eqref{align:ch17-1754}结合，又得
\begin{align}
r'(au+f(u))=s'(au-f(u))=\frac{1}{2a}\hat\eta'(u),
\label{align:ch17-1756}
\end{align}
由此可在 $\mathbb{R}$ 上确定 $r$ 与 $s$，从而在 $\mathbb{R}^2$ 上确定 $\eta$ 与 $q$。特别地，$\hat\eta''\geqslant0$（在 $\mathbb{R}$ 上）蕴含 $r''\geqslant0$、$s''\geqslant0$（在 $\mathbb{R}$ 上），进而 $\eta_{vv}\geqslant0$（在 $\mathbb{R}^2$ 上）。由于 $\eta_v(u,f(u))=0$，我们进而得出耗散性条件\eqref{align:ch05-Zalpha-def}成立：
\begin{align}
\eta_v(u,v)[v-f(u)]\geqslant0,\qquad(u,v)\in\mathbb{R}^2.
\label{align:ch17-1757}
\end{align}
在更强假设 $\hat\eta''(u)\geqslant\beta>0$（$u\in\mathbb{R}$）下，\eqref{align:ch17-1757}变为更严格的
\begin{align}
\eta_v(u,v)[v-f(u)]\geqslant\gamma|v-f(u)|^2,\qquad(u,v)\in\mathbb{R}^2,
\label{align:ch17-1758}
\end{align}
其中 $\gamma>0$。

至此，我们已为确立\eqref{align:ch17-1751}的 Cauchy 问题解的存在性、以及为令 $\mu\downarrow0$ 取极限铺平了道路。

\begin{theorem}[松弛格式的整体存在性]\label{theorem:ch17-松弛存在性}
在次特征条件\eqref{align:ch17-1752}下，方程组\eqref{align:ch17-1751}以
\begin{align}
(u(x,0),v(x,0))=(u_{0\mu}(x),v_{0\mu}(x)),\qquad-\infty<x<\infty,
\label{align:ch17-1759}
\end{align}
为初值（初值含于 $L^\infty(-\infty,\infty)\cap L^2(-\infty,\infty)$）的 Cauchy 问题，在上半平面上有一个有界（弱）解 $(u_\mu,v_\mu)$。此外，
\begin{align}
\frac{1}{\mu}\int_0^\infty\int_{-\infty}^\infty[v_\mu-f(u_\mu)]^2\,dx\,dt\leqslant b\int_{-\infty}^\infty[u_{0\mu}^2(x)+v_{0\mu}^2(x)]\,dx,
\label{align:ch17-17510}
\end{align}
其中 $b$ 与 $\mu$ 无关。
\end{theorem}

\begin{proof}
由于\eqref{align:ch17-1751}是半线性双曲型，局部解 $(u_\mu,v_\mu)$ 存在，并且只要在 $L^\infty$ 中保持有界就可延拓。此外，若 $(\eta,q)$ 是任意熵--熵流对，
\begin{align}
\partial_t\eta(u_\mu,v_\mu)+\partial_x q(u_\mu,v_\mu)+\frac{1}{\mu}\eta_v(u_\mu,v_\mu)[v_\mu-f(u_\mu)]=0.
\label{align:ch17-17511}
\end{align}

我们构造由\eqref{align:ch17-1755}诱导、且由 $\eta_m(0,0)=0$、$q_m(0,0)=0$ 规范化的熵--熵流对 $(\eta_m,q_m)$，其中 $\hat\eta(u)=|u|^m$，$m=2,3,\ldots$。注意 $\eta_m$ 的一阶导数在原点也必为零。将\eqref{align:ch17-17511}在 $(-\infty,\infty)\times[0,t]$ 上积分并使用\eqref{align:ch17-1757}，得
\begin{align}
\int_{-\infty}^\infty\eta_m(u_\mu(x,t),v_\mu(x,t))\,dx\leqslant\int_{-\infty}^\infty\eta_m(u_{0\mu}(x),v_{0\mu}(x))\,dx.
\label{align:ch17-17512}
\end{align}
由\eqref{align:ch17-1756}与\eqref{align:ch17-1752}易知，$(\hat c|w|)^m\leqslant r_m(w)\leqslant(\hat C|w|)^m$ 且 $(\hat c|w|)^m\leqslant s_m(w)\leqslant(\hat C|w|)^m$，从而
\begin{align}
c^m(|u|^m+|v|^m)\leqslant\eta_m(u,v)\leqslant C^m(|u|^m+|v|^m),\qquad(u,v)\in\mathbb{R}^2.
\label{align:ch17-17513}
\end{align}
于是将\eqref{align:ch17-17512}取 $1/m$ 次幂并令 $m\to\infty$，我们断言 $\|u_\mu(\cdot,t)\|_{L^\infty(-\infty,\infty)}$ 与 $\|v_\mu(\cdot,t)\|_{L^\infty(-\infty,\infty)}$ 由 $\|u_{0\mu}(\cdot)\|_{L^\infty(-\infty,\infty)}$ 与 $\|v_{0\mu}(\cdot)\|_{L^\infty(-\infty,\infty)}$ 一致地（关于 $t$ 与 $\mu$）控制。因此解 $(u_\mu,v_\mu)$ 存在于整个上半平面。

接着对熵--熵流对 $(\eta_2,q_2)$ 写出\eqref{align:ch17-17511}，并在 $(-\infty,\infty)\times[0,\infty)$ 上积分。对此情形适用更强的耗散性不等式\eqref{align:ch17-1758}，从而推出\eqref{align:ch17-17510}。证明完成。
\end{proof}

\begin{theorem}[松弛极限]\label{theorem:ch17-松弛极限}
考虑 Cauchy 问题\eqref{align:ch17-1751}、\eqref{align:ch17-1759}的解族 $\{(u_\mu,v_\mu)\}$，其中 $\{(u_{0\mu},v_{0\mu})\}$ 含于 $L^\infty(-\infty,\infty)\cap L^2(-\infty,\infty)$ 的一个有界集，且当 $\mu\downarrow0$ 时 $u_{0\mu}\rightharpoonup u_0$ 在 $L^\infty(-\infty,\infty)$ 弱$^*$ 意义下成立。则存在序列 $\{\mu_j\}$（$\mu_j\downarrow0$，$j\to\infty$），使得 $\{(u_{\mu_j},v_{\mu_j})\}$ 在 $L^\infty$ 弱$^*$ 意义下收敛于 $(\bar u,f(\bar u))$，其中 $\bar u$ 是上半平面上\eqref{align:ch17-1741}的解，初值为 $\bar u(x,0)=u_0(x)$（$-\infty<x<\infty$）。此外，若使 $f''(u)\neq0$ 的 $u$ 之集在 $\mathbb{R}$ 中稠密，则 $\{(u_{\mu_j},v_{\mu_j})\}$ 在上半平面上几乎处处收敛于 $(\bar u,f(\bar u))$。
\end{theorem}

\begin{proof}
由\cref{theorem:ch17-松弛存在性}，$\{(u_\mu,v_\mu)\}$ 含于空间 $L^\infty((-\infty,\infty)\times[0,\infty))$ 的一个有界集。

固定\eqref{align:ch17-1741}的任意熵--熵流对 $(\hat\eta,\hat q)$，考虑由求解 Cauchy 问题\eqref{align:ch17-1753}、\eqref{align:ch17-1755}生成的\eqref{align:ch17-1751}的熵--熵流对 $(\eta,q)$，并利用\eqref{align:ch17-17511}写出
\begin{align}
\begin{split}
\partial_t\hat\eta(u_\mu)+\partial_x\hat q(u_\mu)
={}&\partial_t[\eta(u_\mu,f(u_\mu))-\eta(u_\mu,v_\mu)]+\partial_x[q(u_\mu,f(u_\mu))-q(u_\mu,v_\mu)]\\
&-\frac{1}{\mu}\eta_v(u_\mu,v_\mu)[v_\mu-f(u_\mu)].
\end{split}
\label{align:ch17-17514}
\end{align}

由\eqref{align:ch17-17510}，$\eta(u_\mu,f(u_\mu))-\eta(u_\mu,v_\mu)$ 与 $q(u_\mu,f(u_\mu))-q(u_\mu,v_\mu)$ 当 $\mu\downarrow0$ 时在 $L^2$ 中都趋于零。因此，\eqref{align:ch17-17514}右端前两项当 $\mu\downarrow0$ 时在 $W^{-1,2}$ 中趋于零。另一方面，第三项含于 $L^1$ 的一个有界集，这同样由\eqref{align:ch17-17510}（并回忆 $\eta_v(u,f(u))=0$）得来。

现固定任意序列 $\{\mu_k\}$（$\mu_k\downarrow0$，$k\to\infty$），并令 $(u_k,v_k)=(u_{\mu_k},v_{\mu_k})$。由上述及\cref{lemma:ch17-插值紧致引理}，\eqref{align:ch17-1742}对\eqref{align:ch17-1741}的任意熵--熵流对 $(\hat\eta,\hat q)$ 成立，其中 $\Omega$ 为上半平面。于是\cref{theorem:ch17-标量补偿紧致}给出某子列 $\{u_j\}$ 所满足的\eqref{align:ch17-1743}。进而，\eqref{align:ch17-1743}连同\eqref{align:ch17-17510}蕴含 $v_j\to f(\bar u)$ 在 $L^\infty$ 弱$^*$ 意义下。特别地，由\eqref{align:ch17-1751}$_1$，$\bar u$ 是以 $u_0$ 为初值的\eqref{align:ch17-1741}的解。

当使 $f''(u)\neq0$ 的 $u$ 之集在 $\mathbb{R}$ 中稠密时，由\cref{theorem:ch17-标量补偿紧致}，$\{u_j\}$ 几乎处处收敛于 $\bar u$。再由\eqref{align:ch17-17510}，同理 $\{v_j\}$ 几乎处处收敛于 $f(\bar u)$。证明完成。
\end{proof}

结合\eqref{align:ch17-17511}、\eqref{align:ch17-1757}、\eqref{align:ch17-17510}与\eqref{align:ch17-1755}，我们推断：至少在 $\{u_j\}$ 几乎处处收敛的情形下，极限 $\bar u$ 对任意 $\hat\eta$ 为凸的熵--熵流对 $(\hat\eta,\hat q)$ 满足熵可容性条件。

注意\cref{theorem:ch17-松弛极限}对 $v_\mu$ 的初值 $v_{0\mu}$ 未加任何限制，只要求它们有界。特别地，$v_{0\mu}$ 可偏离其局部平衡值 $f(u_{0\mu})$ 很远。在该情形下，$v_k$ 必在 $t=0$ 附近发展出一个边界层。

应当告诫读者：补偿紧致性并非处理简单方程组\eqref{align:ch17-1751}的最有效方法。事实上，已有结论（文献见\cref{section:ch17-notes}）表明：若 $(u,v)$ 与 $(\bar u,\bar v)$ 是\eqref{align:ch17-1751}的任意一对解，对应初值为 $(u_0,v_0)$ 与 $(\bar u_0,\bar v_0)$，则
\begin{align}
\begin{split}
&\int_{-\ell}^{\ell}\{|u(x,t)-\bar u(x,t)|+|v(x,t)-\bar v(x,t)|\}\,dx\\
&\qquad\leqslant\frac{(1+a)^2}{a}\int_{-\ell-at}^{\ell+at}\{|u_0(x)-\bar u_0(x)|+|v_0(x)-\bar v_0(x)|\}\,dx
\end{split}
\label{align:ch17-17515}
\end{align}
对任意 $\ell>0$、$t>0$ 成立。凭借这一估计，人们可容易地在 $L^1$ 与 $BV$ 中确立紧致性，然后过渡到 $\mu\downarrow0$ 松弛极限。然而，在撰写本书时，对于非线性方程组\eqref{align:ch05-relaxed-system}，
补偿紧致性仍是唯一可行的方法，因为该情形下尚无与\eqref{align:ch17-17515}类似的估计。

\end{document}
